\begin{thebibliography}{31}
\bibitem{small2006growth} H. Small, Tracking and predicting growth areas in science. \textit{Scientometrics} \textbf{68}, 595–610 (2006). doi:10.1007/s11192-006-0132-y
\bibitem{jin2021prizes} C. Jin, Y. Ma, B. Uzzi, Scientific prizes and the extraordinary growth of scientific topics. \textit{Nat. Commun.} \textbf{12}, 5619 (2021). doi:10.1038/s41467-021-25712-2
\bibitem{milojevic2018temporary} S. Milojević, F. Radicchi, J. P. Walsh, Changing demographics of scientific careers: The rise of the temporary workforce. \textit{Proc. Natl. Acad. Sci. U.S.A.} \textbf{115}, 12616–12623 (2018). doi:10.1073/pnas.1800478115
\bibitem{crane1969fashion} D. Crane, Fashion in Science: Does It Exist?. \textit{Soc. Probl.} \textbf{16}, 433–441 (1969). doi:10.2307/799952
\bibitem{fujimura1988bandwagon} J. H. Fujimura, The Molecular Biological Bandwagon in Cancer Research: Where Social Worlds Meet. \textit{Soc. Probl.} \textbf{35}, 261–283 (1988). doi:10.2307/800622
\bibitem{abrahamson1996fashion} E. Abrahamson, Management Fashion. \textit{Acad. Manage. Rev.} \textbf{21}, 254 (1996). doi:10.2307/258636
\bibitem{bikhchandani1992fads} S. Bikhchandani, D. Hirshleifer, I. Welch, A Theory of Fads, Fashion, Custom, and Cultural Change as Informational Cascades. \textit{J. Polit. Econ.} \textbf{100}, 992–1026 (1992). doi:10.1086/261849
\bibitem{banerjee1992herd} A. V. Banerjee, A Simple Model of Herd Behavior. \textit{Q. J. Econ.} \textbf{107}, 797–817 (1992). doi:10.2307/2118364
\bibitem{lorenzspreen2019attention} P. Lorenz-Spreen, B. M. Mønsted, P. Hövel, S. Lehmann, Accelerating dynamics of collective attention. \textit{Nat. Commun.} \textbf{10}, 1759 (2019). doi:10.1038/s41467-019-09311-w
\bibitem{wu2007novelty} F. Wu, B. A. Huberman, Novelty and collective attention. \textit{Proc. Natl. Acad. Sci. U.S.A.} \textbf{104}, 17599–17601 (2007). doi:10.1073/pnas.0704916104
\bibitem{kindleberger2005manias} C. P. Kindleberger, R. Z. Aliber, \textit{Manias, Panics and Crashes: A History of Financial Crises} (Palgrave Macmillan, 2005).
\bibitem{shiller2000exuberance} R. J. Shiller, \textit{Irrational Exuberance} (Princeton University Press, 2000).
\bibitem{pedersen2014bubbles} D. B. Pedersen, V. F. Hendricks, Science Bubbles. \textit{Philos. Technol.} \textbf{27}, 503–518 (2013). doi:10.1007/s13347-013-0142-7
\bibitem{dedehayir2016hype} O. Dedehayir, M. Steinert, The hype cycle model: A review and future directions. \textit{Technol. Forecast. Soc. Change} \textbf{108}, 28–41 (2016). doi:10.1016/j.techfore.2016.04.005
\bibitem{intemann2022hype} K. Intemann, Understanding the Problem of “Hype”: Exaggeration, Values, and Trust in Science. \textit{Can. J. Philos.} \textbf{52}, 279–294 (2020). doi:10.1017/can.2020.45
\bibitem{ioannidis2022covidization} J. P. A. Ioannidis, E. Bendavid, M. Salholz-Hillel, K. W. Boyack, J. Baas, Massive covidization of research citations and the citation elite. \textit{Proc. Natl. Acad. Sci. U.S.A.} \textbf{119}, e2204074119 (2022). doi:10.1073/pnas.2204074119
\bibitem{yeoteh2021covid} N. S. L. Yeo-Teh, B. L. Tang, An alarming retraction rate for scientific publications on Coronavirus Disease 2019 (COVID-19). \textit{Account. Res.} \textbf{28}, 47–53 (2020). doi:10.1080/08989621.2020.1782203
\bibitem{rzhetsky2015choosing} A. Rzhetsky, J. G. Foster, I. T. Foster, J. A. Evans, Choosing experiments to accelerate collective discovery. \textit{Proc. Natl. Acad. Sci. U.S.A.} \textbf{112}, 14569–14574 (2015). doi:10.1073/pnas.1509757112
\bibitem{foster2015tradition} J. G. Foster, A. Rzhetsky, J. A. Evans, Tradition and Innovation in Scientists’ Research Strategies. \textit{Am. Sociol. Rev.} \textbf{80}, 875–908 (2015). doi:10.1177/0003122415601618
\bibitem{jia2017interest} T. Jia, D. Wang, B. K. Szymanski, Quantifying patterns of research-interest evolution. \textit{Nat. Hum. Behav.} \textbf{1}, 0078 (2017). doi:10.1038/s41562-017-0078
\bibitem{hill2025pivot} R. Hill, Y. Yin, C. Stein, X. Wang, D. Wang, B. F. Jones, The pivot penalty in research. \textit{Nature} \textbf{642}, 999–1006 (2025). doi:10.1038/s41586-025-09048-1
\bibitem{fortunato2018science} S. Fortunato et al., Science of science. \textit{Science} \textbf{359}, eaao0185 (2018). doi:10.1126/science.aao0185
\bibitem{bornmann2015growth} L. Bornmann, R. Mutz, Growth rates of modern science: A bibliometric analysis based on the number of publications and cited references. \textit{J. Assoc. Inf. Sci. Technol.} \textbf{66}, 2215–2222 (2015). doi:10.1002/asi.23329
\bibitem{priem2022openalex} J. Priem, H. Piwowar, R. Orr, OpenAlex: A fully-open index of scholarly works, authors, venues, institutions, and concepts. arXiv:2205.01833 (preprint, 2022); \url{https://doi.org/10.48550/arXiv.2205.01833}.
\bibitem{waltman2012classification} L. Waltman, N. J. van Eck, A new methodology for constructing a publication‐level classification system of science. \textit{J. Am. Soc. Inf. Sci. Technol.} \textbf{63}, 2378–2392 (2012). doi:10.1002/asi.22748
\bibitem{johansen2000crashes} A. Johansen, O. Ledoit, D. Sornette, Crashes as critical points. \textit{Int. J. Theor. Appl. Finance} \textbf{3}, 219–255 (2000). doi:10.1142/S0219024900000115
\bibitem{sornette2003stock} D. Sornette, \textit{Why Stock Markets Crash: Critical Events in Complex Financial Systems} (Princeton University Press, 2003).
\bibitem{salganik2006experimental} M. J. Salganik, P. S. Dodds, D. J. Watts, Experimental Study of Inequality and Unpredictability in an Artificial Cultural Market. \textit{Science} \textbf{311}, 854–856 (2006). doi:10.1126/science.1121066
\bibitem{ke2015sleeping} Q. Ke, E. Ferrara, F. Radicchi, A. Flammini, Defining and identifying Sleeping Beauties in science. \textit{Proc. Natl. Acad. Sci. U.S.A.} \textbf{112}, 7426–7431 (2015). doi:10.1073/pnas.1424329112
\bibitem{retractionwatch2023} The Center for Scientific Integrity, The Retraction Watch Database. Crossref and Retraction Watch, https://api.labs.crossref.org/data/retractionwatch (accessed September 2026) (2026).
\bibitem{brainard2018retractions} J. Brainard, J. You, What a massive database of retracted papers reveals about science publishing's `death penalty'. Science, News, 25 October 2018 (2018); doi:10.1126/science.aav8384.
\end{thebibliography}
